%=====================================================================
%  BACHELOR THESIS  (version 1: Introduction and Literature Review)
%
%  Confidence-Aware Adaptive Verification of LLM Recommendations
%  in Introductory Programming Tasks
%
%  Student:    Lidia Davydova
%  Supervisor: Paolo Ciancarini
%
%  This file is a self-contained first draft. It contains the front
%  matter, the Introduction (Chapter 1), and the Literature Review
%  (Chapter 2). Later chapters are placeholders to be filled in.
%=====================================================================

\documentclass[oneside,final,14pt,a4paper]{extreport}

\usepackage{tempora} % Times New Roman alike font

\usepackage{vmargin}
\setpapersize{A4}
\setmarginsrb{2.5cm}{2.2cm}{2.2cm}{2.2cm}{0pt}{10mm}{0pt}{13mm}
\usepackage{setspace}
\setstretch{1.5}
\usepackage{indentfirst}
\parindent=1.25cm

% Fix the bold monospaced face and use Courier for \texttt.
\DeclareFontShape{OT1}{cmtt}{bx}{n}{<5><6><7><8><9><10><10.95><12><14.4><17.28><20.74><24.88>cmttb10}{}
\renewcommand{\ttdefault}{pcr}

\usepackage[T1]{fontenc}
\usepackage[utf8]{inputenc}
\usepackage[english]{babel}

\usepackage{amsmath,amsfonts,amssymb}
\usepackage{amsthm}
\usepackage{upgreek}
\usepackage{bm}

\usepackage{graphicx}
\graphicspath{{figs/}}
\usepackage{float}
\usepackage{array}
\usepackage{multirow}
\usepackage{booktabs}
\usepackage{longtable}
\usepackage[font=singlespacing,labelfont=bf]{caption}
\usepackage{subcaption}
\usepackage{enumitem}
\usepackage{csquotes}
\usepackage{url}
\usepackage{xcolor}

\usepackage[backend=biber,style=ieee,autocite=inline]{biblatex}
\addbibresource{references.bib}
\DefineBibliographyStrings{english}{%
  bibliography = {References},}

\usepackage{chngcntr}
\usepackage{hyperref}
\hypersetup{colorlinks=true, linkcolor=black, citecolor=black, urlcolor=black}

\usepackage{fancyhdr}
\pagestyle{fancyplain}
\renewcommand{\chaptermark}[1]{\markboth{\chaptername~\thechapter~--~#1}{}}
\renewcommand{\sectionmark}[1]{\markright{\thesection\ #1}}
\rhead[\fancyplain{}{\bf\leftmark}]{\fancyplain{}{\bf\thepage}}
\lhead[\fancyplain{}{\bf\thepage}]{\fancyplain{}{\bf\rightmark}}
\cfoot{}
\setlength{\headheight}{32.0976pt}

\captionsetup[table]{labelfont={normalfont}, name={TABLE}, labelsep={newline}}
\captionsetup[figure]{labelfont={normalfont}, name={Fig.}, labelsep=period}


\begin{document}

%=====================================================================
%  TITLE PAGE
%=====================================================================
\begin{titlepage}
\thispagestyle{empty}
\begin{center}
\vspace*{2cm}
{\large Innopolis University\par}
\vspace{2cm}
{\LARGE\bfseries
Confidence-Aware Adaptive Verification of LLM Recommendations\\
in Introductory Programming Tasks\par}
\vspace{2.5cm}
{\large Student: Lidia Davydova\par}
\vspace{0.5cm}
{\large Supervisor: Paolo Ciancarini\par}
\vfill
{\large \today\par}
\end{center}
\end{titlepage}

\tableofcontents
\newpage

%=====================================================================
%  ABSTRACT
%=====================================================================
\chapter*{Abstract}
\addcontentsline{toc}{chapter}{Abstract}

Large language models (LLMs) are increasingly used in programming education,
but their recommendations can be plausible yet incorrect, so the quality of
human--AI collaboration depends on \emph{appropriate reliance}: accepting
correct recommendations and rejecting incorrect ones. Verification support can
help users evaluate recommendations, and introductory programming is well
suited to it because correctness can be checked through independent,
execution-based evidence such as tests, traces, and counterexamples. Providing
such evidence on every trial, however, introduces unnecessary interaction cost.
This thesis investigates whether a simple, interpretable adaptive policy---gating
verification availability on the participant's initial confidence---can reduce
verification effort relative to always-available verification while preserving
decision quality. The human--AI interaction is specified through an explicit
\emph{escalation rule} that fixes who recommends, who supplies independent
verification evidence, who decides, and what is recorded; the three
experimental conditions are three escalation policies within that single
specification. A within-subject
experiment compares no-verification, static, and adaptive conditions on a
benchmark of introductory Python tasks with frozen, pre-validated LLM
recommendations balanced between correct and incorrect. The two primary
outcomes are final-answer accuracy and verification effort; reliance measures
are secondary. Reliance on the LLM is distinguished from agreement with an
answer the participant already held, and adaptive verification is tested for
non-inferiority in decision quality against static verification, including the
boundary case of confident but mistaken participants. Expected contributions
include trial-level evidence on the effort--quality trade-off of
confidence-gated verification, a validated reproducible benchmark with an
instrumented prototype, and a complete analysis pipeline.

\newpage

%=====================================================================
%  CHAPTER 1 --- INTRODUCTION
%=====================================================================
\chapter{Introduction}
\label{chap:intro}

\section{Motivation}
\label{sec:motivation}

Large language models (LLMs) are increasingly used in programming education and
software engineering to generate, explain, and debug code and to propose
answers to programming tasks \cite{Ma2024,Cambaz2024,Denny2024}. They can lower
the cost of producing a first attempt, suggest alternative implementations, and
explain unfamiliar concepts. At the same time, such systems can produce
plausible but incorrect outputs: fluent code that fails to compile, passes some
tests but not others, or satisfies the letter of a specification while missing
its intent.

Because of this, the quality of the joint human--AI outcome depends not only on
the accuracy of the model but on how the human partner responds to it. This is
captured by the notion of \emph{appropriate reliance} \cite{LeeSee2004}:
relying on correct recommendations and refraining from relying on incorrect
ones, judged relative to the answer the user already holds
\cite{Schemmer2023Appropriate}. Its two failure modes are well documented.
\emph{Over-reliance} occurs when a user who was initially correct accepts an
incorrect recommendation; \emph{under-reliance} occurs when a user who was
initially incorrect rejects a correct one. Both can persist even when additional
information about the model is made available
\cite{PoursabziSangdeh2021,Vasconcelos2022,Bo2025}.

Verification support is a natural remedy: rather than asking users to trust or
distrust the model on the basis of fluency, it lets them check a recommendation
against independent evidence. Programming tasks suit this approach particularly
well, because correctness can often be checked independently through execution,
tests, traces, or counterexamples---evidence grounded in the task itself rather
than in another LLM-generated justification. This distinction matters, because
explanations help only when they sufficiently reduce the effort of checking the
recommendation, and simply providing more information does not automatically
improve appropriate reliance
\cite{Bucinca2021,Vasconcelos2022,Bo2025,Schemmer2022}.
Moreover, feature-based explanations reduce over-reliance mainly for decisions
that users can already evaluate themselves and lose their benefit as difficulty
rises
\cite{Zhang2024Verify}; introductory programming provides exactly such checkable
grounding.

The remaining difficulty is cost. Verification on every trial introduces
unnecessary interaction effort: if evidence is displayed when it is not needed,
it slows users down and consumes attention that could be spent elsewhere. This
motivates providing support selectively: escalating verification effort when
the system's confidence in the current answer is low \cite{CascadedLM2025}, or
when a user-side signal such as trust indicates that support is needed
\cite{SrinivasanAdjustForTrust}, using only information that is available before
the correctness of the recommendation is known.

\section{Research Problem}
\label{sec:problem}

The research problem is how to provide verification support selectively, using
only information available before the correctness of the LLM recommendation is
known. This thesis investigates whether a simple adaptive policy, based
primarily on the user's initial confidence, can provide independent verification
support at lower effort than always-available verification, and how verification
effort trades off against decision quality.

The focus is a simple, interpretable policy---not a new LLM, a new verification
algorithm, or an optimisation method. The contribution is therefore not the
threshold rule itself; it is the controlled empirical evaluation of whether a
simple confidence gate can reduce verification exposure while preserving
appropriate reliance compared with always-available and no-verification
baselines.

The problem is important because verification cost is paid on every trial, while
its benefit is concentrated on the trials where the user is uncertain or wrong.
Introductory programming amplifies both sides: novices are likely to over-rely
on fluent but incorrect code, yet they are also the users least able to afford
repeated checking. A policy that spends verification effort only where it is
needed would make independent evidence practical in ordinary learning settings,
not only in controlled studies.

\section{Research Gap}
\label{sec:gap}

Prior work has examined confidence mainly as an outcome to be calibrated, as an
AI-side uncertainty signal for routing a case to stronger processing
\cite{CascadedLM2025}, or as a user-side trust signal for adapting explanations
\cite{SrinivasanAdjustForTrust}. It has also shown that explanations and
cognitive forcing can reduce over-reliance
\cite{Bucinca2021,Bo2025,Schemmer2022}. What is not known is whether the
participant's own initial confidence can serve as a per-trial trigger for
independent, execution-based verification, and whether such a gate preserves
decision quality while reducing verification effort. Studies that offer support
on every trial do not measure the cost of unnecessary verification, while
studies that manipulate AI-side signals do not test a user-side gate. Direct
evidence on this effort--quality trade-off is therefore missing, especially for
novices in introductory programming, where confident errors are common and
independent evidence is cheap to produce.

The gap matters because a confidence-gated policy would be a lightweight,
model-agnostic way to reduce interruption cost without weakening decisions. If
the policy is not non-inferior, the result identifies when confidence should not
be used as a gate, which is equally important for safe deployment. In
introductory programming the stakes are concrete: novices are prone to
over-reliance on plausible but incorrect code, and uncorrected errors can
entrench misconceptions. The question therefore has to be answered empirically
before such a policy is recommended in practice.

\section{Conceptual Framework: An Explicit Escalation Rule}
\label{subsec:protocol}

Coordination, in the account of Ciancarini \cite{Ciancarini2027}, concerns who
decides, with what evidence, and under whose authority, so that correct and
incorrect recommendations can be managed by the surrounding system rather than
by conversation alone. This thesis adopts that element in an explicit
\emph{escalation rule}: a specification of who supplies evidence, who decides,
and when the evidence reaches the decider. The rule does not add any
implementation; it makes the responsibilities already present in the design
explicit. It has five parts:

\begin{itemize}
    \item \textbf{Recommender.} A fixed, offline LLM snapshot produces the
    candidate answer for each task. The recommender has no access to the
    experimental condition, the verification policy, or the participant.
    \item \textbf{Evidence provider.} The researchers supply independent,
    execution-based verification artefacts (tests, traces, counterexamples)
    derived from the task specification and reference answer, not from the
    recommender. The evidence provider never states whether the recommendation
    is correct.
    \item \textbf{Decider.} The participant holds decision authority: they
    record an initial answer, a confidence rating, and a final judgment.
    \item \textbf{Escalation rule.} The rule specifies when the evidence
    provider's output is escalated to the decider. Decision authority is never
    reassigned: the participant decides in every condition, and only the flow
    of evidence is conditional. Under the adaptive policy the escalation
    condition is $c<\tau$ on the participant's initial confidence; under the
    static policy evidence is always escalated; under no-verification it never
    is.
    \item \textbf{Audit and record.} The instrumented prototype records the
    task, condition, escalation decision, evidence availability and examination,
    timestamps, and the correctness label. The correctness label is kept only in
    the analysis dataset and is never shown to the decider.
\end{itemize}

The three experimental conditions are therefore \emph{three escalation
policies within the same specification}: only the escalation rule is changed,
while the recommender, the evidence provider, the decider, and the audit record
are held fixed. This keeps the implementation unchanged and makes
the comparison a comparison of escalation policies rather than of different
systems.

\section{Research Questions}
\label{sec:rq}

\textbf{Primary RQ.} Does confidence-aware adaptive verification reduce
verification effort relative to always-available (static) verification in
introductory programming tasks, while preserving decision quality relative to
static verification (non-inferiority) and final-answer accuracy relative to no
verification?

\textbf{RQ1 (primary).} How do no, static, and adaptive verification affect
final-answer accuracy?

\textbf{RQ2 (primary).} How does the verification effort of adaptive
verification (verification availability, evidence-examination rate,
evidence-examination time) compare with static verification, and is decision
quality preserved?

\textbf{RQ3 (exploratory).} How do the effort and effectiveness of the adaptive
policy differ across initial confidence levels, task types, and the
confident-but-mistaken stratum?

The primary hypotheses follow from the research questions.

\begin{itemize}
\item \textbf{H1 (quality).} Adaptive verification is non-inferior to static
verification on final-answer accuracy within the pre-specified margin
$\delta_Q$.
\item \textbf{H2 (effort).} Adaptive verification reduces total decision time
relative to static verification.
\end{itemize}

Reliance measures are reported as secondary outcomes, distinguishing true
reliance on the LLM from agreement with an answer the participant already held.
Neither task type nor task difficulty is an input to the adaptive policy; both
are analysed only as exploratory moderators of policy effectiveness. A policy
premise check examines whether initial confidence is associated with initial
answer correctness, and a non-inferiority test establishes whether adaptive
verification preserves decision quality relative to static verification.

\section{Approach}
\label{sec:approach}

The thesis uses a within-subject experiment in which each participant completes
eighteen introductory Python tasks under three escalation policies: no
verification, static verification, and confidence-gated adaptive verification.
Condition order is counterbalanced. Recommendations are produced offline by a
frozen LLM snapshot and pre-validated so that each condition contains a balanced
mix of correct and incorrect advice; verification artefacts are derived from the
task specification and executed independently. The two primary outcomes are
final-answer accuracy and total decision time, with reliance measures reported
as secondary outcomes. Trial-level mixed-effects models account for the nesting
of trials within participants and tasks, and accuracy is assessed by a
non-inferiority test while effort is assessed by a superiority test. The
instrumented prototype records evidence availability and examination so that the
policy can be reproduced and audited.

\section{Contributions and Scope}
\label{sec:contributions}

The novelty is \emph{not} the observation that confidence relates to reliance on
AI---miscalibration is well documented \cite{KrugerDunning1999,Bo2025} and taken
here as a premise---but the combination of elements usually studied in
isolation: (i)~confidence as a per-trial \emph{trigger} deciding \emph{when}
support is shown, rather than an outcome to be calibrated or an AI-side
uncertainty signal; (ii)~\emph{independent, execution-based} evidence (tests,
traces, counterexamples) rather than additional AI-generated explanations;
(iii)~a programming setting in which recommendations are pre-validated and
objective evidence exists; and (iv)~an explicit effort--quality trade-off
evaluation, treating verification effort as a first-class outcome---in contrast
to prior work that measures accuracy, reliance, or system cost
\cite{CascadedLM2025,Zhang2024Verify}.

The thesis is expected to contribute: (1)~trial-level empirical evidence,
analysed with mixed-effects models, on whether a simple confidence-based policy
reduces verification effort while preserving decision quality (RQ1--RQ2);
(2)~an explicit escalation rule that states who recommends, who provides
evidence, who decides, and what is recorded, and that treats the three
conditions as three escalation policies within one specification
\cite{Ciancarini2027}; (3)~a validated benchmark of introductory Python tasks
with fixed, pre-validated LLM recommendations (balanced correct/incorrect),
task-type and difficulty labels, and verification data; (4)~a reproducible
prototype of the three conditions with complete logging, including evidence
examination; and (5)~an analysis pipeline for the primary outcome and effort
measures.

The contribution is deliberately not a new LLM, a new algorithm, or a tutoring
system: the confidence gate is intentionally simple, and its value lies in the
controlled empirical comparison against always-available and no-verification
baselines in a reproducible human--AI interaction experiment.

Taken together, these contributions show how a minimal user-side signal can
govern when independent evidence is worth its cost. If the policy is supported,
it offers programming environments a model-agnostic way to reduce interruption
without weakening decisions; if it is not, the boundary conditions delimit when
confidence gating is unsafe. Either outcome informs the design of verification
support for novice programmers.

\section{Structure of the Thesis}
\label{sec:structure}

The remainder of the thesis is organised as follows.
Chapter~2 reviews the literature on trust and appropriate reliance,
explanation-based, cognitive-forcing, and selective interventions, verification
and its cost, confidence as a signal for adaptive support, the coordination and
governance of LLM-based systems, and LLM assistance in programming education, and
derives the research gap addressed here.
Chapter~3 presents the methodology: the benchmark and task domain, the
recommendation specification, the verification artefacts, the three escalation
policies, and the outcome measures.
Chapter~4 describes the implementation of the instrumented prototype and its
logging.
Chapter~5 reports the evaluation, the statistical analysis, the discussion of
results, threats to validity, and implications.
Chapter~6 concludes and outlines future work.


%=====================================================================
%  CHAPTER 2 --- LITERATURE REVIEW
%=====================================================================
\chapter{Literature Review}
\label{chap:lr}

\section{Introduction and Review Scope}
\label{sec:lit-intro}

Research on human--AI decision-making spans human factors, explainable AI,
human--computer interaction, and computing education. Reviewing these
literatures separately would obscure the tension motivating this thesis: an
intervention that improves trust in one setting can degrade the joint decision
in another, so the relevant question is not whether users trust a system but
whether they rely on it appropriately. The review is organised around four
strands---trust and appropriate reliance; explanation-based, cognitive-forcing,
and adaptive interventions; verification, confidence, and coordination; and
LLM assistance in programming education---and for each it classifies the
approaches, evaluates the evidence, and states the implications for this
study. Sources came from targeted, iterative searches of the major computing
and HCI databases, prioritising peer-reviewed work from 2019--2027 and
including older foundational work where it defines a construct used here
\cite{LeeSee2004,KrugerDunning1999}. The review is argument-driven rather than
systematic, and its final set is restricted to studies that report an
empirical human--AI interaction result, define a construct, metric, or design
framework used in the thesis, or describe a comparable educational setting.

\section{Trust, Reliance, and Calibration}
\label{sec:lit-trust}

Trust in automation is classically defined as the attitude that an agent will
help an individual achieve a goal under uncertainty and vulnerability
\cite{LeeSee2004}. Lee and See's contribution is not only the definition but
the design goal that follows from it: the objective is not maximal trust but
\emph{appropriate} reliance, calibrated to the actual capability of the
automated agent. Appropriate reliance concerns the extent to which a decision
maker relies on correct advice and refrains from relying on incorrect advice
\cite{Schemmer2023Appropriate}. Its two failure modes are defined relative to
the initial answer. \emph{Over-reliance} occurs when a decision maker who was
initially correct follows incorrect advice; \emph{under-reliance} occurs when a
decision maker who was initially incorrect rejects correct advice
\cite{Schemmer2023Appropriate,Bo2025}. This conditioning is essential: without
it, accepting wrong advice and rejecting correct advice cannot be distinguished
from cases in which the initial answer was already wrong or already correct.

The distinction matters because interventions that reduce over-reliance may 
inadvertently increase under-reliance, potentially without improving overall 
reliance appropriateness \cite{Bo2025}. The construct is not merely terminological: how it is
operationalized determines which errors a study can detect. Schemmer et al.
observe that appropriate reliance still lacks a common definition and an
operational measurement concept; their two-dimensional formulation, adopted
here, separates reliance on correct advice from reliance on incorrect advice
and reinforces the need to state explicitly which errors a policy is meant to
reduce.

Empirical work in AI-assisted decision-making has shown that users frequently fail
to calibrate reliance to recommendation quality. In a sequence of pre-registered
experiments with functionally identical models that varied only in apparent
interpretability, participants did not follow predictions more closely when the
model seemed clearer, and in some conditions were \emph{less} able to detect
the model's mistakes, apparently because of information overload
\cite{PoursabziSangdeh2021}. Chen et al. \cite{Chen2023} traced this result to
the interaction between explanations and users' own intuition about the
outcome. Feature-based explanations were visually dominant: participants
scanned highlighted spans rather than reading, more often reported weak
intuition about the outcome, and deferred to the AI when they could not reason
about the highlighted features, so decision outcomes did not improve and
over-reliance increased. Example-based explanations left the instance and the
user's prior judgment in the foreground, so strong intuition could be used to
override the AI and the explanation could be ignored, supporting complementary
performance.
Bo et al. \cite{Bo2025} benchmarked three interventions with LLM advice in a
randomised experiment with 400 participants: a static reliance disclaimer
appended to the advice, token-level uncertainty highlighting, and an implicit
answer that withheld the recommendation and required participants to derive it.
All three reduced over-reliance, but only the disclaimer significantly improved
appropriate reliance, and only on the LSAT task; in numerical estimation no
intervention did, and participants became more confident after incorrect
reliance decisions in some contexts. These studies establish that over- and
under-reliance are persistent, not incidental, and that adding information
does not by itself produce calibration.

A second body of work concerns the relation between confidence and
correctness. Kruger and Dunning \cite{KrugerDunning1999} reported that less
competent individuals tend to overestimate their own performance relative to
more competent peers, a pattern usually described as a metacognitive deficit.
The finding is often invoked as evidence that self-assessment is unreliable,
but the underlying phenomenon is broader and more differentiated: Moore and
Healy \cite{MooreHealy2008} distinguish overestimation of one's absolute
performance, overplacement relative to others, and overprecision in the
certainty of one's beliefs, and show that these components do not always move
together. For the present thesis, the relevant implication is narrower than
the popular Dunning--Kruger narrative: confidence is informative about
correctness on average, but it is not a direct readout of correctness, and its
diagnostic value must be checked empirically in the task domain rather than
assumed.

These strands jointly define the outcome vocabulary used in this thesis. At
the trial level, \emph{appropriate reliance} occurs when the participant was
initially incorrect, received correct advice, and adopted it, or was initially
correct, received incorrect advice, and rejected it. \emph{Over-reliance} is
the proportion of trials on which the participant was initially correct,
received incorrect advice, and adopted it; \emph{under-reliance} is the
proportion of trials on which the participant was initially incorrect, received
correct advice, and rejected it. Conditioning on the initial answer also
separates reliance from mere agreement: when the participant already holds the
recommended answer, agreement with the LLM is not evidence of reliance on it.
This is why the thesis distinguishes \emph{agreement} ($F_i=L_i$, the final
answer matches the recommendation) from \emph{LLM influence} ($A_i\neq L_i$ and
$F_i=L_i$, the participant moved toward the recommendation from a different
prior answer) and treats stable prior agreement as a separate case. The
distinction is developed further in the chapter's discussion of evaluation
metrics.

\section{Interventions That Support Reliance}
\label{sec:lit-explanations}

Research on supporting reliance can be classified by the mechanism it targets:
the information shown to the user, the interaction that forces deliberation,
the expectations with which the user enters the task, or the policy that decides
when support is provided. These families are not mutually exclusive, but the
distinction is useful because they rest on different assumptions about why
reliance fails and impose different costs on the user.

\subsection{Explanations and Explainable AI}
\label{sec:lit-xai}

Explanations are the most extensively studied intervention. The rationale is
that revealing why a model produced a recommendation should help users detect
errors and calibrate their trust. The empirical record is more conditional. A
meta-analysis of nine articles covering 44 studies and a range of explanation
types---feature importance, example-based explanations, and confidence
scores---found that XAI-assisted decisions outperformed unaided human decisions
by 0.59 standard deviations (95\% CI 0.39--0.79). In the narrower subset of
studies that compared explanations with an AI recommendation alone, however,
the difference was only 0.07 (95\% CI $-0.15$--$0.30$): the pooled estimate
concealed substantial heterogeneity, with a prediction interval ranging from
$-0.59$ to $0.74$, so explanations helped in some settings and hurt in others
\cite{Schemmer2022}. The gain over unaided humans was concentrated
in text tasks (0.72) and absent for tabular ones (0.01). A qualitative reading
of the same studies concluded that explanations mainly raised acceptance of AI
advice, with only one article reporting evidence of appropriate reliance. This
distinction matters: explanations may make an AI system more acceptable without
making its advice easier to evaluate.

Vasconcelos et al. \cite{Vasconcelos2022} offer a cost--benefit account that
reconciles some of the mixed results. Across five studies with 731
participants, they manipulated task difficulty, explanation difficulty, and
monetary incentives, and found that explanations reduce over-reliance when they
sufficiently lower the cost of verifying the prediction. When checking an
explanation is as costly as solving the task, users rationally fall back on the
recommendation. The benefit of an explanation therefore lies in the verification
effort it saves, not in its presence.

Explanation format also matters. Chen et al. \cite{Chen2023} compared
feature-based and example-based explanations in a think-aloud study and found
that feature-based explanations increased over-reliance, whereas example-based
explanations supported complementary human--AI performance. Their account
centres on whether an explanation displaces or scaffolds the user's existing
intuition about the outcome: feature-based displays disrupted it, whereas
example-based displays preserved it and supplied independent signals of AI
unreliability. Zhang et al. \cite{Zhang2024Verify} sharpened this point by
manipulating decision difficulty directly. Feature-based explanations reduced
over-reliance for easy decisions but not for hard ones, and for the hardest
decisions they could increase it.
Their title states the underlying constraint: users can only verify when they
already know the answer. Bansal et al. \cite{Bansal2021} reached a compatible
conclusion from the perspective of team performance: explanations increased the
chance that participants accepted the AI's recommendation regardless of whether
it was correct, and did not increase complementary team performance. Wang et
al. \cite{Wang2019} make the broader design argument: an explanation should be
chosen for the reasoning goal it serves and the cognitive bias it is meant to
counter, rather than from a fixed list of explanation types or from whatever
XAI technique happens to be available.

Taken together, these studies show that explanations are neither necessary nor
sufficient for appropriate reliance. They can help when they reduce the cost of
checking a claim the user is able to check, and they can hurt when they persuade
through fluency rather than evidence. This is one reason the present thesis
separates explanation from verification: an explanation is generated by or
about the model, whereas verification evidence is obtained from the task
itself, independently of the model.

\subsection{Cognitive Forcing and Expectation Setting}
\label{sec:lit-forcing}

A second family changes the interaction rather than the content. Cognitive
forcing functions change when and how a recommendation becomes available, so
that accepting it requires an act of engagement rather than a passive read. The
three designs tested by Bu\c{c}inca et al. \cite{Bucinca2021} illustrate the
range: participants had to request the AI's suggestion explicitly, to commit to
an initial answer before seeing it and then revise, or to wait 30 seconds while
the system appeared to process the image. Compared with simple explainable-AI
interfaces and a no-AI baseline in an experiment with 199 participants, these
interventions significantly reduced over-reliance but did not eliminate it, and
the designs that reduced over-reliance most received the least favourable
subjective ratings. The benefit was also unevenly distributed: participants higher in Need
for Cognition benefited more, which the authors report as an
intervention-generated inequality. Cognitive forcing therefore trades
effectiveness against acceptability and may widen differences between users.

Expectation setting acts earlier, before the interaction begins. Kocielnik et
al. \cite{Kocielnik2019} studied a scheduling assistant whose AI component had
the same 50\% accuracy in two conditions that differed only in whether the
design emphasised avoiding false positives or false negatives. These framings
produced markedly different subjective perceptions of accuracy and acceptance,
and explicit expectation-adjustment techniques increased acceptance of the
imperfect system. This work is relevant because it shows that users' evaluation
of an imperfect system is shaped by framing and prior expectations, not only by
its accuracy. Its outcome, however, is acceptance rather than decision accuracy:
baseline and expectation-setting conditions did not differ significantly in
objective task accuracy (0.91 vs.\ 0.88), so preparing users for imperfection
can make a system more acceptable without making their decisions better.

Both approaches change behaviour, but neither supplies task-grounded evidence
that the user can inspect. Cognitive forcing can interrupt automatic acceptance,
and expectation setting can calibrate the user's prior, yet both leave the
burden of verification on the user and can reduce reliance overall rather than
direct it toward correct recommendations.

\subsection{Adaptive and Selective Interventions}
\label{sec:lit-adaptive}

A third family makes support conditional. The condition may be a property of
the model, of the task, or of the user. Fanconi and van der Schaar
\cite{CascadedLM2025} propose a cascaded LLM framework in which a base model
answers first, a deferral policy sends uncertain cases to a larger model, and an
abstention policy sends remaining uncertain cases to a human expert; the
thresholds are updated online from human feedback. The framework balances
accuracy, computational cost, and human effort, and it is evaluated on
question-answering benchmarks rather than in a controlled human-subject study.
Its deferral signal is the model's calibrated confidence, while abstention to
the human is triggered by posterior uncertainty; and the human receives a case
to decide, not a body of evidence to evaluate.

The work closest to this thesis on the user side is that of Srinivasan and
Thomason \cite{SrinivasanAdjustForTrust}. They argue that extreme trust produces
under-reliance when trust is low and over-reliance when it is high, and adapt
the assistant's behaviour accordingly: supporting explanations are shown when
trust is low, counter-explanations when trust is high, and forced pauses are
inserted to promote deliberation. In two decision-making scenarios, including
medical diagnosis, supporting and counter-explanations reduced inappropriate
reliance by up to 38\% and improved decision accuracy by up to 20\%; forced
pauses separately reduced over-reliance by promoting deliberation. The signal
is the user's trust state, and the gated support is still AI-generated
explanation or interaction pacing, not independent verification of the
recommendation.

The intervention studies reviewed above are mostly static: the same support is
provided on every trial. Bo et al. \cite{Bo2025} evaluate three such
interventions with LLM advice and find that they reduce over-reliance but
generally fail to improve appropriate reliance, because they reduce reliance
overall rather than redirect it. The authors discuss confidence calibration as
a design implication but do not use confidence as a control signal. This
motivates the present work, which combines a user-side, per-trial confidence
signal with independent execution-based evidence.

\subsection{Comparative Evaluation}
\label{sec:lit-comparison}

Table~\ref{tab:positioning} compares the closest lines of work along four
dimensions: the signal that drives support, what is gated, what the user sees,
and the primary outcome. The comparison makes the gap concrete. Model-side
gating uses confidence to decide whether to escalate to a more capable model or
to a human \cite{CascadedLM2025}; trust-adaptive gating uses a user state to
choose an explanation or a pause \cite{SrinivasanAdjustForTrust}; the
intervention studies apply their support uniformly \cite{Bo2025}; and Zhang et
al.\ vary decision difficulty but do not gate on it, finding that feature
explanations reduce over-reliance only for easy decisions
\cite{Zhang2024Verify}. Within the scope of the search described in
Section~\ref{sec:lit-intro}, no study has been found that gates \emph{independent,
execution-based} evidence on the decider's own prior confidence, and none has
treated interaction effort as a primary outcome alongside decision quality.

\begin{table}[t]
\centering
\caption{Positioning relative to the closest lines of work.}
\label{tab:positioning}
\footnotesize
\begin{tabular}{@{}>{\raggedright\arraybackslash}p{2.4cm}>{\raggedright\arraybackslash}p{2.0cm}>{\raggedright\arraybackslash}p{2.6cm}>{\raggedright\arraybackslash}p{3.2cm}>{\raggedright\arraybackslash}p{2.6cm}@{}}
\toprule
Work & Signal & What is gated & What the user sees & Primary outcome \\
\midrule
Bo et al. \cite{Bo2025} & None (fixed) & Nothing (uniform) & Disclaimer, uncertainty highlighting, or implicit answer & Appropriate reliance \\
Srinivasan and Thomason \cite{SrinivasanAdjustForTrust} & User trust state & Explanation type; forced pause & Supporting or counter-explanations & Inappropriate reliance \\
Zhang et al. \cite{Zhang2024Verify} & None (difficulty as moderator) & Nothing (explanation vs.\ no explanation) & Feature attributions & Over-reliance \\
Fanconi and van der Schaar \cite{CascadedLM2025} & Calibrated confidence; posterior uncertainty & Large-model escalation; human review & Case only (no evidence) & Cost and accuracy \\
This thesis & Prior-answer confidence & Independent execution evidence & Tests, traces, counterexamples & Accuracy and effort \\
\bottomrule
\end{tabular}
\end{table}

The comparison also exposes a methodological difference. The studies above
measure decision quality, reliance, or computational cost; CascadedLM measures
human involvement as a cost to be minimised, but not the interaction effort of
the person who receives the case. The present thesis treats verification effort
as a first-class outcome because a policy that preserves accuracy by making
everyone examine evidence on every trial has not solved the problem it set out
to solve.

\section{Independent Verification and Its Cost}
\label{sec:lit-verification}

\emph{Verification} here means checking a claim against evidence that is
independent of the source that produced it. In programming, the evidence can be
a test suite, an execution trace, or a counterexample derived from the task
specification and the reference answer. This differs from an explanation in
kind, not only in presentation. An explanation is an account of why the model
produced an answer; verification evidence is a check on whether the answer
satisfies the task. The distinction is important because a fluent explanation
can increase acceptance without increasing correctness \cite{Bansal2021}, and
because an explanation of a model's reasoning need not faithfully describe the
computation that produced the answer.

Verification is not free. Human scrutiny of AI outputs consumes time and
attention, and the cost is incurred on every trial on which evidence is shown.
Vasconcelos et al. \cite{Vasconcelos2022} formalise this as a cost--benefit
choice: users weigh the cognitive effort of verifying the AI against the
lower-effort strategy of relying on it, and engage with an explanation when the
reduction in verification cost justifies the effort. CascadedLM
\cite{CascadedLM2025} applies a related cost logic at the system level, routing a
case to human review only after a larger model still returns a high-uncertainty
answer. Zhang et al. \cite{Zhang2024Verify} add a constraint that matters for
novices in particular: verification presupposes an independent reference against
which the answer can be judged. Feature-based explanations reduced overreliance
only on decisions participants could already evaluate themselves, and not on
difficult ones. Evidence that cannot be checked within the decision-time budget
is not effective support, however objective it is.

Introductory programming is a favourable setting for independent verification.
Tasks come with executable test suites, so correctness can be checked by
running tests rather than by judging surface similarity; code-generation
benchmarks adopt functional correctness for exactly this reason
\cite{Liu2023EvalPlus}. The tasks are small enough that a trace or a failing test
can plausibly be interpreted within the decision-time budget, and the evidence
can be derived from the task specification rather than from the recommender. The
evidence is still not exhaustive: a program can pass a finite test suite and
remain incorrect, and a trace shows what happened on one input rather than on
all inputs. Execution-based
evidence therefore supports judgment rather than replacing it, which is
consistent with leaving decision authority with the participant.

The design problem is consequently not how to construct verification evidence
but \emph{when} to present it. Showing evidence on every trial pays the cost
even when the recommendation is already accepted correctly; withholding it
entirely leaves the participant with fluency as the only cue. The next section
considers whether the participant's own confidence can serve as the signal that
decides when the cost is worth paying.

\section{Confidence as a Signal for Adaptive Support}
\label{sec:lit-confidence}

Lee and See \cite[70]{LeeSee2004} observe that static approaches used to
identify miscalibration of self-confidence in decision making ``address only
cumulative experience, rather than the evolving experience and continuous
recalibration that are critical for appropriate reliance on automation.'' The
implication drawn here is that calibration should not be treated as a stable
property to be estimated once and then applied as a fixed weight; what matters is
the relation between confidence and correctness as it evolves across tasks,
trials, and moments of an interaction.

This distinction matters for adaptive verification. A single calibration score,
or a disposition measured before the task, cannot distinguish a trial on which a
student is correctly confident from one on which the same student is confidently
wrong. A per-trial confidence signal can, in principle, make that distinction at
the point where support is being considered. Throughout this thesis, confidence
refers to the participant's metacognitive judgement about the correctness of
their own prior answer, elicited before the recommendation is shown; it is
distinct both from trust in the assistant \cite{LeeSee2004} and from the model's
own uncertainty. The thesis therefore adopts the weaker, testable implication of
Lee and See's argument: confidence is elicited on every trial and used as a
routing signal, while its relationship to initial-answer correctness is checked
empirically rather than assumed. The stronger claim---that the decision
threshold itself should be continuously recalibrated---is outside the scope of
the present design.

Empirically, the premise is plausible but not guaranteed. Bo et al.
\cite{Bo2025} found that interventions reduced over-reliance but generally
failed to improve appropriate reliance, and that confidence after incorrect
reliance decisions was poorly calibrated in some contexts. The authors discuss
confidence calibration as a design implication but do not use confidence as a
control signal. Related work adapts support to the user's trust state
\cite{SrinivasanAdjustForTrust}, adjusting the interaction to the user's
evolving disposition toward the system.

In operational terms, confidence is an \emph{input signal}---a lightweight
trigger deciding \emph{when} independent support is shown---rather than a fully
calibrated internal state. This is a deliberately modest operationalisation of
the dynamic view: confidence is measured repeatedly and used at the trial level,
but the threshold is fixed rather than learned online. The documented
relationship between confidence and correctness is treated as a premise: if
confidence carries information about correctness, it can also inform \emph{when}
verification is worth its interaction cost. The interaction-cost dimension is
largely absent from calibration-oriented work and is a primary outcome here.

\section{Coordination and Governance of LLM-Based Systems}
\label{sec:lit-coordination}

Because a language model can produce a fluent answer without guaranteeing that
it is correct, in scope, or compatible with other decisions, plausible but
unsupported output must be managed by the surrounding system rather than trusted
on fluency alone \cite{Ciancarini2027}. Recent work on agentic software
engineering argues that increasing autonomy raises, rather than lowers, the need
for coordination: claims may need to be checked against artefacts, actions
against permissions, and decisions against explicitly assigned authority, and
shared artefacts, logs, tests, and approval points make decisions traceable
\cite{Ciancarini2027}.

This perspective motivates the explicit escalation rule adopted in this thesis
(Section~\ref{subsec:protocol}): coordination concerns who decides, with what
evidence, and under whose authority, and here that allocation is made explicit
in independent, execution-based evidence that replaces reliance on
conversational output alone.
It also clarifies why verification is not merely a UX feature: it is the
mechanism by which the system, rather than the fluency of the model, governs
when a recommendation is trusted.

\section{LLMs in Programming Education}
\label{sec:lit-programming}

LLMs are increasingly used for code generation, debugging, and explanation in
computing education \cite{Ma2024,Cambaz2024,Denny2024}. They can accelerate
routine work and provide on-demand feedback, but they can also be plausible yet
incorrect, so learners must evaluate recommendations critically. Prior work
establishes that intuition, explanation design, and decision difficulty shape
reliance \cite{Chen2023,Vasconcelos2022,Zhang2024Verify}. In educational
settings, the stakes of miscalibration are high: accepting an incorrect
recommendation can entrench a misconception, while rejecting a correct one
wastes an opportunity to learn.

\section{Synthesis and Research Gap}
\label{sec:lit-gap}

The reviewed literature converges on three points. First, appropriate
reliance---not trust maximisation---is the correct objective, and over- and
under-reliance are persistent failure modes \cite{LeeSee2004,PoursabziSangdeh2021,
Bo2025}. Second, explanations and cognitive forcing help only conditionally,
and information alone does not guarantee better monitoring
\cite{Bucinca2021,Schemmer2022,Vasconcelos2022}. Third, human scrutiny is costly
and should be invoked selectively \cite{CascadedLM2025}, and explanation-based
support helps mainly on decisions users can already evaluate themselves
\cite{Zhang2024Verify}.

What remains unknown is whether a learner's \emph{initial confidence} can serve
as a lightweight trigger for selectively displaying \emph{independent
verification evidence} after an LLM recommendation in introductory programming
tasks---reducing verification effort relative to always-available support while
preserving final-answer accuracy and appropriate reliance. This is the gap the
thesis addresses.

Since LLM-based assistants have entered programming education, the cost of
human scrutiny has become a primary design constraint, and the reviewed
literature points to a controlled comparison of escalation policies rather
than to a new model or algorithm. Accordingly, the policy is tested in a
within-subject experiment in which all participants experience all three
escalation policies, so that trial-level comparisons are not confounded by
individual differences. Trial-level outcomes are analysed with mixed-effects
models because trials are nested within participants and tasks. Decision
quality is assessed with a non-inferiority test rather than a superiority
test, because the cost literature reviewed here indicates that the goal is to
preserve quality while reducing verification effort. Recommendations are
frozen and pre-validated so that the verification policy, not model behaviour,
is the manipulated factor. Concretely, the design treats confidence as an
escalation signal within a single explicit escalation rule, measures
verification effort as a first-class outcome, and evaluates decision quality
against explicit non-inferiority criteria, including the boundary case of
participants who are confident but mistaken.


%=====================================================================
%% REFERENCES
\printbibliography[heading=bibintoc,title={References}]

\end{document}
